Ens trobem en un indret en què el camp magnètic terrestre té una magnitud de $35\ \mathrm{\mu T}$ i apunta cap al nord, però està inclinat 50° (cap avall) respecte a l'horitzontal (figura $a$). Per un cable d'una línia d'alta tensió situada en aquest indret hi circula un corrent de 200~A d'intensitat. Aquest corrent circula d'est a oest (cap endins a la figura $a$, $\otimes$).

\figura[0.42]{fig1.png}

\begin{apartat}{1,25}
Calculeu la força magnètica que actua sobre un tram de 100 metres del cable deguda al camp magnètic terrestre. Després, determineu-ne el mòdul i representeu-ne esquemàticament la direcció i el sentit. Justifiqueu la resposta.
\end{apartat}

\begin{apartat}{1,25}
El corrent que circula pel cable d'alta tensió és un corrent altern i genera un camp magnètic que contínuament canvia de sentit (vegeu la figura $b$).

\figura[0.42]{fig2.png}

A sota del cable s'han situat tres espires conductores: una (A) és paral·lela a la superfície horitzontal del terreny, una altra (B) és paral·lela al pla del dibuix, i la tercera (C) està situada en el pla vertical que conté la direcció est-oest. En quina o quines de les espires el camp magnètic variable produït per la línia d'alta tensió induirà un corrent elèctric? Justifiqueu la resposta especificant les lleis o els principis físics en què us heu basat.
\end{apartat}

\nota{Considereu que el camp magnètic és uniforme en la regió on es troben les espires.}
